% Dataset section. Paste into your own thesis, or \input it.

\section{Coronary tree structure}
\label{sec:dataset-tree}

The left and right coronary systems arise from separate ostia (Section~\ref{sec:tree}), and
ImageCAS-X stores them as separate structures, one centerline file per side per case
\cite{bransby2026imagecasx}.

In the 789 cohort cases where it is labeled, the left main trunk has a median length of 8.5 mm
(interquartile range 6.0--11.6 mm, full range 1.7--30.4 mm), so the first and largest branching
point of the left tree sits within a centimeter of the aorta.

\begin{figure}[h]
  \centering
  \includegraphics[width=\textwidth]{content/background/figures/anatomy_ct.png}
  \caption{The left main divides within millimeters of the aorta, and the annotation locates it to
    the voxel. Axial computed tomography (CT) slice of case 7 at the level of the coronary ostia,
    cutting the aortic root rather than the ventricles, with the delivered segmentation overlaid in
    the per-artery colors used throughout. The ostium marker is the centerline's own
    \texttt{start\_points} flag converted into voxel indices, where it samples 627 Hounsfield units
    (HU), the density scale of a CT image. Generated with
    \texttt{scripts/make\_anatomy\_figure.py}.}
  \label{fig:anatomy-ct}
\end{figure}

Prevalence of the numbered diagonal and obtuse marginal branches (Section~\ref{sec:tree}) falls
with position: D1 is labeled in 98.5\% of cases and D2 in 59.4\%, OM1 in 86.8\% and OM2 in 37.6\%.

The right system is annotated more coarsely.
ImageCAS-X's label set represents only the RCA trunk and the posterior vessels, so the proximal
right-sided branches described in Section~\ref{sec:tree} are not labeled
\cite{bransby2026imagecasx}.
This is an annotation boundary rather than an anatomical claim, and it constrains what a
comparison of branch counts between the left and right trees can mean.
